\documentclass[aps,prd,twocolumn,amsmath,amssymb,floatfix,nofootinbib]{revtex4-2}

\usepackage[T2A]{fontenc}
\usepackage[utf8]{inputenc}
\usepackage[english,russian]{babel}
\usepackage{graphicx}
\usepackage{bm}
\usepackage{hyperref}
\hypersetup{colorlinks=true,linkcolor=blue,citecolor=blue,urlcolor=blue}

\begin{document}

\title{Точное численное соотношение, связывающее отношение масс
тау-лептона и электрона с электромагнитной и гравитационной
константами связи}

\author{Thomas J.\ Buckholtz}
\affiliation{National Coalition of Independent Scholars, Brattleboro, Vermont 05301, USA}
\affiliation{Ronin Institute for Independent Scholarship, Montclair, New Jersey 07043, USA}
\email{tjb@alumni.caltech.edu}

\date{\today}

\begin{abstract}
Мы сообщаем о численном соотношении между измеряемыми величинами,
$(4/3)\,(m_\tau/m_e)^{12} = \alpha_{\mathrm{EM}}/\alpha_G$,
где $\alpha_{\mathrm{EM}}$~--- постоянная тонкой структуры, а
$\alpha_G \equiv G m_e^2/(\hbar c)$~--- гравитационная константа связи,
отнесённая к массе электрона. При использовании текущих данных Particle
Data Group и CODATA обе части согласуются с точностью $0{,}061\%$, что
соответствует $1{,}0\,\sigma$ современной погрешности массы тау-лептона.
Чтобы оценить, могло ли такое согласие возникнуть случайно, мы перебираем
все соотношения вида $(p/q)(m_i/m_j)^n$ при $p,q\le 10$, $1\le n\le 24$,
где $m_i,m_j$ выбираются из одиннадцати масс заряженных фермионов, протона
и нейтрона; среди полученных $83\,160$ кандидатов данное соотношение
занимает первое место, а следующий по точности кандидат менее точен в
$3{,}4$ раза. Если позволить опорной массе также меняться среди всех
одиннадцати кандидатов, ничего не меняется: соотношение~(\ref{eq:main})
остаётся единственным кандидатом первого ранга, уже с отрывом в $6{,}8$
раза. Эмпирическая вероятность случайного совпадения на уровне одного
процента ниже $10^{-4}$. Мы приводим соотношение, его статистическую
значимость, алгебраическую структуру его целочисленных коэффициентов, а
также остающиеся открытыми вопросы~--- прежде всего вывод коэффициента
$4/3$, который, как мы обнаруживаем, не даётся ни одной смежной
теоретической традицией. Мы полагаем, что соотношение относится, в
статической (не зависящей от времени) форме, к традиции совпадений больших
чисел, включающих гравитационную константу связи.
\end{abstract}

\maketitle

\section{Введение}

Соотношения между массами элементарных частиц, а также между этими массами
и фундаментальными константами связи, издавна используются как указатели на
пути к более глубокой теории~\cite{koide1983,dirac1938,wesson1980}. Наиболее
известное такое соотношение среди заряженных лептонов~--- формула
Койде~\cite{koide1983}, которая включает только три массы заряженных
лептонов и не содержит отсылки к гравитации. Отдельно гипотеза больших чисел
П.~А.~М.~Дирака~\cite{dirac1938} обратила внимание на огромное безразмерное
отношение между электромагнитной и гравитационной константами связи и на его
кажущееся совпадение с космологическими величинами.

Здесь мы сообщаем о соотношении, связывающем эти две темы: оно связывает
отношение масс заряженных лептонов с отношением электромагнитной и
гравитационной констант связи. Мы формулируем соотношение, оцениваем,
насколько маловероятно, что оно случайно, исследуем арифметику его
коэффициентов и~--- в интересах полной прозрачности~--- прямо указываем, что
соотношение устанавливает, а что нет. Наша цель~--- зафиксировать один
проверяемый численный факт, а не заявить механизм для него.

\section{Соотношение}

Определим гравитационную константу связи, отнесённую к массе электрона, в
духе Дирака~\cite{dirac1938},
\begin{equation}
\alpha_G \equiv \frac{G\,m_e^2}{\hbar c},
\label{eq:alphaG}
\end{equation}
которая является гравитационным аналогом постоянной тонкой структуры
$\alpha_{\mathrm{EM}} = e^2/(4\pi\varepsilon_0\hbar c)$. Выбор массы
электрона в качестве опорной~--- это соглашение; мы возвращаемся к нему в
разд.~\ref{sec:discussion}.

Сообщаемое нами соотношение имеет вид
\begin{equation}
\frac{4}{3}\left(\frac{m_\tau}{m_e}\right)^{\!12}
= \frac{\alpha_{\mathrm{EM}}}{\alpha_G}.
\label{eq:main}
\end{equation}
С входными данными табл.~\ref{tab:inputs} обе части вычисляются как
\begin{align}
\text{ЛЧ} &= 4{,}16814\times 10^{42}, \nonumber\\
\text{ПЧ} &= 4{,}16561\times 10^{42},
\end{align}
так что
\begin{equation}
\frac{\text{ЛЧ}}{\text{ПЧ}} = 1{,}000608,
\end{equation}
относительное расхождение $0{,}0608\%$. Поскольку левая часть растёт как
$m_\tau^{12}$, текущая погрешность массы тау-лептона
$\sigma_{m_\tau}=0{,}09~\mathrm{МэВ}$ переносится в
одностандартную полосу $12\,\sigma_{m_\tau}/m_\tau = 0{,}061\%$ на левой
части. Таким образом, наблюдаемое расхождение составляет $1{,}0\,\sigma$~---
обычное согласие масштаба одного измерения, а не подозрительно точное.
Эквивалентно, соотношение~(\ref{eq:main}) выполняется точно при
\begin{equation}
m_\tau = 1776{,}840~\mathrm{МэВ},
\end{equation}
что отстоит на $1{,}0\,\sigma$ от текущего центрального значения PDG,
$1776{,}93\pm 0{,}09~\mathrm{МэВ}$~\cite{pdg2024}. Измерение массы
тау-лептона с точностью $0{,}01~\mathrm{МэВ}$, ожидаемой на
Belle~II~\cite{belleii}, усилит проверку соотношения~(\ref{eq:main}) на
порядок, до $\text{ЛЧ}/\text{ПЧ} = 1{,}0000 \pm 0{,}0001$. Отметим, что это
подлинно двусторонняя проверка, а не заведомое подтверждение: при
пересчёте на точность Belle~II \emph{текущее} центральное значение PDG
$m_\tau=1776{,}93$~МэВ отстоит на $9{,}0\,\sigma_{\text{Belle II}}$ от
соотношения~(\ref{eq:main})~--- измерение Belle~II, согласующееся с
сегодняшним центральным значением на ожидаемой точности, решительно
опровергло бы точное соотношение, тогда как сдвиг к $1776{,}840$~МэВ его
подтвердил бы. То, что нынешняя масса~--- измеренная с точностью порядка
$0{,}1\%$~--- уже отстоит лишь на $1{,}0\,\sigma$ (а не, скажем, на
$5\,\sigma$), и есть актуальное сегодня утверждение о совместимости;
цифра $9\,\sigma$ описывает лишь то, насколько остро сможет различать
Belle~II в будущем.

\begin{table}[t]
\caption{Входные величины и их источники. Массы в
$\mathrm{МэВ}/c^2$.}
\label{tab:inputs}
\begin{ruledtabular}
\begin{tabular}{lll}
Величина & Значение & Источник \\
\hline
$m_e$ & $0{,}51099895000$ & CODATA~2018~\cite{codata2018} \\
$m_\tau$ & $1776{,}93\pm 0{,}09$ & PDG~2024~\cite{pdg2024} \\
$\alpha_{\mathrm{EM}}$ & $1/137{,}035999084$ & CODATA~2018~\cite{codata2018} \\
$G$ & $6{,}67430\times 10^{-11}$ & CODATA~2018~\cite{codata2018} \\
$\alpha_G$ & $1{,}75181\times 10^{-45}$ & из ур.~(\ref{eq:alphaG}) \\
\end{tabular}
\end{ruledtabular}
\end{table}

\section{Статистическая значимость}
\label{sec:lookelsewhere}

Одно близкое численное совпадение само по себе ничего не доказывает: имея
достаточную свободу в выборе масс, показателей и рациональных множителей,
можно приблизить почти любую целевую величину. Поэтому существенен вопрос о
поправке на множественность (look-elsewhere)~--- насколько особым является
соотношение~(\ref{eq:main}) в семействе соотношений, которые можно было бы
записать вместо него?

Чтобы ответить, мы перечисляем каждое соотношение вида
\begin{equation}
\frac{p}{q}\left(\frac{m_i}{m_j}\right)^{\!n},
\qquad
1\le p,q \le 10,\quad
1\le n \le 24,\quad
i\neq j,
\label{eq:family}
\end{equation}
где $m_i,m_j$ выбираются из множества
$\{e,\mu,\tau,u,d,s,c,b,t,p,n\}$, и сравниваем каждое с целевой величиной
$\alpha_{\mathrm{EM}}/\alpha_G$. Это задаёт пространство поиска из $83\,160$
кандидатов. Результаты сведены в табл.~\ref{tab:lookelsewhere}.

\begin{table}[t]
\caption{Аудит множественности соотношения~(\ref{eq:main}) по семейству
из $83\,160$ членов~(\ref{eq:family}). «Попадания»~--- число кандидатов,
согласующихся с целевой величиной точнее указанного допуска.}
\label{tab:lookelsewhere}
\begin{ruledtabular}
\begin{tabular}{lr}
Величина & Значение \\
\hline
Всего кандидатов & $83\,160$ \\
Попаданий в пределах $5\%$ & $25$ \\
Попаданий в пределах $1\%$ & $5$ \\
Попаданий в пределах $0{,}1\%$ & $1$ \\
Ранг ур.~(\ref{eq:main}) & $1$ \\
Погрешность ур.~(\ref{eq:main}) & $0{,}0608\%$ \\
Второй по точности & $(2/5)(m_t/m_n)^{19}$, $0{,}21\%$ \\
Эмпирическое $p$ (в пределах $1\%$) & $<10^{-4}$ \\
\end{tabular}
\end{ruledtabular}
\end{table}

Соотношение~(\ref{eq:main})~--- единственный кандидат, согласующийся с
целевой величиной точнее $0{,}1\%$; оно занимает первое место в целом, а
следующее по точности соотношение $(2/5)(m_t/m_n)^{19}$ менее точно в $3{,}4$
раза. Лишь пять из $83\,160$ кандидатов согласуются точнее $1\%$, что даёт
эмпирическую вероятность ниже $10^{-4}$ того, что совпадение такого качества
возникнет случайно в этом семействе. Подчеркнём, что этот анализ
количественно оценивает неслучайность; сам по себе он не даёт механизма.

\subsection{Устойчивость к выбору опорной массы}
\label{sec:refmass}

Ещё один вопрос: масса $m_e$ выбрана опорной в $\alpha_G$ по соглашению
(разд.~\ref{sec:discussion}); не возникло бы столь же точное соотношение при
другой опорной массе, и не обязана ли тогда кажущаяся точность
соотношения~(\ref{eq:main}) размеру пространства поиска, а не электрону?
Поэтому мы повторяем перебор~(\ref{eq:family}), позволяя самой опорной массе
в $\alpha_G$ меняться по всем одиннадцати массам. Просуммированные по всем
одиннадцати вариантам опорной массы, каждая из $55$ неупорядоченных пар масс
перебирается ровно один раз (с более лёгким членом пары в роли естественной
опорной массы), так что объединённое семейство вновь состоит из $83\,160$
кандидатов~--- того же размера, что и в разд.~\ref{sec:lookelsewhere}, но
теперь с плавающей, а не фиксированной опорной массой. Соотношение~
(\ref{eq:main}) остаётся единственным кандидатом первого ранга: ни одна
другая комбинация опорной массы, пары масс, рационального множителя и
показателя не согласуется со своей целью лучше. Следующий по точности
кандидат во всём этом расширенном пространстве,
\begin{equation}
\frac{4}{3}\left(\frac{m_b}{m_s}\right)^{\!23}
= \frac{\alpha_{\mathrm{EM}}}{\alpha_G(m_s)},
\end{equation}
согласуется лишь с точностью $0{,}42\%$~--- хуже соотношения~(\ref{eq:main})
в $6{,}8$ раза, а не точнее. Совокупная эмпирическая вероятность~--- доля
кандидатов из $83\,160$, не хуже соотношения~(\ref{eq:main})~--- составляет
$1/83\,160=1{,}2\times 10^{-5}$: снятие соглашения об опорной массе ничуть
не ослабляет результат, а \emph{усиливает} его, поскольку
соотношение~(\ref{eq:main}) остаётся единственным выбросом в поиске вдвое
большего размера, чем тот, который его впервые выявил.

Более ранняя версия этого анализа, основанная на с тех пор уточнённой массе
топ-кварка, указывала на кандидата
$(4/7)(m_t/m_\tau)^{18}=\alpha_{\mathrm{EM}}/\alpha_G(m_\tau)$, который,
казалось, согласуется со своей целью с точностью $0{,}0070\%$~--- лучше,
чем~(\ref{eq:main}). При текущей массе топ-кварка PDG~2024,
$M_t=172{,}57$~ГэВ (прямые измерения), ошибка этого кандидата составляет
$2{,}03\%$; при $\overline{\mathrm{MS}}$-массе PDG~2024,
$\overline{m}_t(\overline{m}_t)=162{,}5$~ГэВ~\cite{pdg2024}, она составляет
$66{,}8\%$. При любой схеме он теперь далеко за пределами окна точнее $1\%$
и не входит даже в первую десятку объединённого поиска. Мы фиксируем это
как предостережение относительно 18-й степени у этого кандидата: столь
высокий показатель усиливает сдвиг в несколько ГэВ между последовательными
определениями массы топ-кварка PDG до изменения предсказанного отношения
порядка единицы, так что точное совпадение, построенное на такой степени,
следует считать предварительным, пока входная масса не установлена с
точностью лучше процента~--- а для топ-кварка это не так даже между двумя
определениями из одного и того же издания PDG. У соотношения~
(\ref{eq:main}) такой уязвимости нет: $m_e$ и $m_\tau$~--- on-shell
лептонные массы без сопоставимой схемной свободы, а оба его перебора
множественности (с фиксированной и с плавающей опорной массой) устойчивы
именно к той поправке, которая устранила прежнего конкурента.

\section{Гипотетическая алгебраическая интерпретация коэффициентов}
\label{sec:algebra}

Целые числа, входящие в соотношение~(\ref{eq:main})~--- множитель $4/3$ и
показатель $12$~--- допускают наводящие алгебраические прочтения, которые мы
приводим как наблюдения, а не как выводы. Мы не утверждаем, что
$\mathrm{G}_2$, $\mathrm{F}_4$ или исключительная йорданова алгебра
\emph{вызывают} соотношение~(\ref{eq:main}); мы лишь отмечаем, что если
будущий вывод коэффициентов существует, это естественные объекты, которые
он мог бы задействовать.

Показатель $12$ совпадает с двумя независимыми инвариантами исключительных
алгебр Ли: это число ненулевых корней
$\mathrm{G}_2 = \mathrm{Aut}(\mathbb{O})$, группы автоморфизмов октонионов,
и это наибольшая из четырёх степеней Казимира $(2,6,8,12)$ группы
$\mathrm{F}_4$~\cite{bincer1993}. Множитель можно записать как
$4/3 = 36/27$, где $27 = \dim J_3(\mathbb{O})$~--- размерность
исключительной йордановой алгебры (алгебры Альберта), а
$36 = \dim \mathrm{SO}(9)$. Перебор рациональных комбинаций размерностей
классических и исключительных алгебр Ли находит и отношение $4/3$, и
конкретную тройку $\{36,27,12\}$ с оценочной вероятностью случайности
порядка $1/37000$.

Мы явно помечаем этот алгебраический поиск как \emph{вторичный}, найденный
постфактум тест на множественность: в отличие от основного перебора из
разд.~\ref{sec:lookelsewhere}, который был определён и выполнен без ссылки
на значение $4/3$, алгебраический поиск проводился \emph{после} того, как
$4/3$ уже было известно, по набору размерностей, отобранных именно потому,
что ожидалось найти в нём это число. Поэтому величину $1/37000$ следует
читать как меру того, насколько особенна тройка $\{36,27,12\}$
\emph{внутри этого набора}, а не как независимое подтверждение
соотношения~(\ref{eq:main}) на равных с разд.~\ref{sec:lookelsewhere}.

Эти связи указывают алгебраические объекты, которые \emph{подсчитывают}
коэффициенты, но не объясняют, \emph{почему} именно эти объекты управляют
отношением $\alpha_{\mathrm{EM}}/\alpha_G$. В частности, нам не удалось
вывести коэффициент $4/3$ из первых принципов и~--- как обсуждается в
разд.~\ref{sec:discussion}~--- мы не нашли такого вывода ни в одной смежной
теоретической традиции.

\section{Связь с предшествующими работами}

Соотношение~(\ref{eq:main}) отлично от соотношений заряженных лептонов типа
Койде--Сумино~\cite{koide1983,sumino2009} и от программы отношений масс на
исключительной йордановой алгебре~\cite{singh2021}. Те построения
предсказывают отношения между массами фермионов из одной лишь флейворной
симметрии или алгебры и не содержат гравитационной величины; данное же
соотношение характеризуется именно появлением гравитационной константы связи
$\alpha_G$. Мы проверили, что число $12$ входит в алгебру
работы~\cite{singh2021} как счёт собственных значений, а не как показатель
степени, так что численное совпадение этих двух вхождений не означает общего
вывода.

Естественный контекст для соотношения между комбинацией масс частиц и
$\alpha_G$~--- это скорее традиция совпадений больших
чисел~\cite{dirac1938,wesson1980}. В отличие от исходной гипотезы Дирака,
соотношение~(\ref{eq:main}) \emph{статично}: оно связывает величины
современной эпохи и не подразумевает, что $G$ меняется с космологическим
временем. Поэтому оно не наследует наблюдательных трудностей, ограничивающих
переменную $G$.

\section{Обсуждение: что это устанавливает, а что нет}
\label{sec:discussion}

Мы считаем существенными следующие разграничения.

\emph{Что мы сообщаем.} Численное соотношение между измеряемыми величинами,
ур.~(\ref{eq:main}), которое выполняется с точностью $0{,}061\%$
($1{,}0\,\sigma$) и является единственным членом семейства из $83\,160$
кандидатов с точностью лучше $0{,}1\%$~--- и остаётся им, с ещё большим
отрывом, при плавающей опорной массе в $\alpha_G$~--- с эмпирической
вероятностью случайности ниже $10^{-4}$ в обоих переборах.

\emph{Чего мы не утверждаем.} Мы не утверждаем вывода
соотношения~(\ref{eq:main}) из лежащей в основе динамики и не заявляем
механизма, фиксирующего его коэффициенты. Множитель $4/3$ и показатель $12$
здесь не выводятся. Мы обследовали три смежные традиции, в которых можно было
бы ожидать происхождения коэффициента $4/3$~--- флейворную симметрию SU(3),
геометрическое прочтение формулы Койде и традицию больших чисел~--- и ни в
одной из них $4/3$ не возникает как выведенная величина. Поэтому мы
представляем происхождение коэффициента как по-настоящему открытую проблему,
по трудности сравнимую с четырёхдесятилетней проблемой вывода формулы Койде.

\emph{О выборе $m_e$.} Опорная масса в $\alpha_G$ условна. Её роль здесь
самосогласована, а не фундаментальна: одна и та же масса электрона входит и
в отношение $m_\tau/m_e$, и в $\alpha_G$, так что обе части
соотношения~(\ref{eq:main}) выражены в электронных единицах.

\emph{Связь с более широкой программой.} Соотношение~(\ref{eq:main})
предлагается здесь не впервые: в эквивалентной форме
$(4/3)(m_\tau^2/m_e^2)^6 = (1/4\pi\varepsilon_0)q_e^2/(Gm_e^2)$ оно фигурирует
как ур.~(32) более широкого препринта MULTING/IDM, из которого оно
взято~\cite{buckholtzv6}. Там же, с явной осторожностью, предлагается
феноменологическое обоснование показателя степени: <<The exponent 6 in
Eq.~(32) might associate with the number, six, of isomers>> (степень 6 в
ур.~(32) может быть связана с числом шесть изомеров), где шесть~---
число изомеров элементарных частиц (один изомер обычной материи и пять
изомеров тёмной материи), постулируемых в этой программе. Мы приводим это
обоснование для полноты и сохраняем его исходную осторожность формулировки;
мы не подтверждаем и не проверяем саму изомерную программу здесь, и это
обоснование не зависит от алгебраических наблюдений
разд.~\ref{sec:algebra}.

\section{Воспроизводимость}

Все приведённые здесь величины вычисляются из опубликованных констант
табл.~\ref{tab:inputs} элементарной арифметикой. Перечисление множественности
в разд.~\ref{sec:lookelsewhere}~--- детерминированный перебор; скрипт,
полностью воспроизводящий табл.~\ref{tab:lookelsewhere} вместе с вычислением
ур.~(\ref{eq:main}), доступен у автора по запросу и предназначен для
публичного размещения вместе с настоящей заметкой.

\section{Заключение}

Мы зафиксировали одно проверяемое соотношение,
$(4/3)(m_\tau/m_e)^{12} = \alpha_{\mathrm{EM}}/\alpha_G$, которое согласуется
с измерением с точностью $0{,}061\%$ ($1{,}0\,\sigma$), является единственным
выбросом широкого семейства сравнимых соотношений как при фиксированной,
так и при плавающей опорной массе (разд.~\ref{sec:refmass}), и
коэффициенты которого допускают прочтение через исключительные алгебры,
происхождение которого ещё предстоит объяснить. Отражает ли соотношение
глубокую связь между иерархией масс
лептонов и гравитационной константой связи или же пока необъяснённое
совпадение~--- вопрос, который мы оставляем открытым; предстоящее измерение
массы тау-лептона на Belle~II проверит его решающим образом.

\begin{acknowledgments}
Автор благодарит С.~Бойко за проведение независимого анализа множественности
разд.~\ref{sec:lookelsewhere} и за критическое прочтение рукописи.
\end{acknowledgments}

\begin{thebibliography}{99}

\bibitem{koide1983} Y.~Koide,
\textit{Phys.\ Lett.\ B} \textbf{120}, 161 (1983).

\bibitem{dirac1938} P.~A.~M.~Dirac,
\textit{Proc.\ R.\ Soc.\ Lond.\ A} \textbf{165}, 199 (1938).

\bibitem{wesson1980} P.~S.~Wesson,
\textit{Gravity, Particles, and Astrophysics} (Reidel, Dordrecht, 1980).

\bibitem{pdg2024} S.~Navas \textit{et al.} (Particle Data Group),
\textit{Phys.\ Rev.\ D} \textbf{110}, 030001 (2024).

\bibitem{codata2018} E.~Tiesinga, P.~J.~Mohr, D.~B.~Newell, and B.~N.~Taylor,
\textit{Rev.\ Mod.\ Phys.} \textbf{93}, 025010 (2021) [CODATA~2018].

\bibitem{belleii} W.~Altmannshofer \textit{et al.} (Belle~II),
\textit{Prog.\ Theor.\ Exp.\ Phys.} \textbf{2019}, 123C01 (2019).

\bibitem{bincer1993} A.~M.~Bincer and K.~Riesselmann,
\textit{J.\ Math.\ Phys.} \textbf{34}, 5935 (1993).

\bibitem{sumino2009} Y.~Sumino,
\textit{J.\ High Energy Phys.} \textbf{05} (2009) 075.

\bibitem{singh2021} V.~Bhatta, R.~Mondal, V.~Vaibhav, and T.~P.~Singh,
arXiv:2108.05787 [hep-ph] (2021).

\bibitem{buckholtzv6} T.~J.~Buckholtz, <<Gravitational and Dark-Matter
Concepts that Can Help Explain and Predict Cosmic Data>>, Preprints.org,
doi:10.20944/preprints202511.0598.v6 (2026), ур.~(32) и разд.~4.7.

\end{thebibliography}

\end{document}
